\documentclass{ahuthesis}

% 封面元信息
\thesistitlecn1{安徽大学理工类本科毕业论文模板}
\thesistitlecn2{最小示例}
\thesistitleen{A Minimal Thesis Template Example for Anhui University}

\studentname{张三}
\studentid{2022000000}
\department{物理学院}
\major{应用物理}
\enrolldate{2022 年 9 月}
\supervisor{李四}
\supervisortitle{教授 / 博士}
\supervisorunit{安徽大学物理学院}
\completiondate{2026 年 6 月}
\keywordsCN{论文模板；最小示例；排版}
\keywordsEN{thesis template; minimal example; typesetting}

\begin{document}

% 前置部分：封面、摘要、目录
\frontmatter
\makecover

% 中文摘要
\begin{cnabstract}
本文档是安徽大学理工类本科毕业论文模板的最小示例，保留了封面、摘要、目录、正文、参考文献、附录与致谢等基本结构。实际写作时，通常只需修改本文件中的元信息与正文内容，再按需补充更多章节。
\end{cnabstract}

% 英文摘要
\begin{enabstract}
This document is a minimal example of the Anhui University science thesis template. It keeps the basic structure of a thesis, including the cover, abstracts, table of contents, main text, bibliography, appendix, and acknowledgements.
\end{enabstract}

% 生成目录
\tableofcontents

% 正文示例
\mainmatter
\section{引言}
本文档仅保留最基本的示例内容，用于说明论文模板的组织方式。使用时，可以先替换题目、作者信息和正文，再根据需要继续扩展章节、小节与图表等内容。

\subsection{技术路线：使用建议}
为避免与真实出版物混淆，下面的引文与文后参考文献均采用纯虚构示例，只用于展示不同 BibTeX 条目在模板中的排版效果。条目中的“作者甲”“出版社甲”“地点甲”“Note Alpha”等均是字段占位符，用来明确字段归属。

图书类与版次信息可参见 \cite{demoBookCN,demoBookEN}；期刊类的卷(期)与页码格式可参见 \cite{demoArticleCN,demoArticleEN}。
会议论文、技术报告与学位论文的展示效果可参见 \cite{demoProceedingCN,demoReportEN,demoThesisCN}。实际写作时，只需按相同字段结构替换为自己的真实文献信息即可。

\section{结论}
通过保持示例内容精简，模板可以更清楚地展示论文的基本结构，也便于作者直接替换为自己的正式内容。

% 文后部分：参考文献、附录、致谢
\backmatter
\bibliographystyle{ahu_citation}
\bibliography{references}

% 附录
\ahuappendix{A}{附录示例}
此处可放置补充材料，例如原始数据、推导过程或额外说明。

% 致谢
\begin{acknowledgements}
感谢导师、同学与家人在论文写作过程中提供的帮助与支持。
\end{acknowledgements}

\end{document}
